\chapter{Polyarteritis Nodosa}
\index{Polyarteritis Nodosa}
\label{sec:Polyarteritis Nodosa}

\section{Overview}
Polyarteritis nodosa is a systemic necrotizing vasculitis affecting medium-sized arteries, characterized by inflammation and fibrinoid tissue death of the vessel wall, with aneurysm formation and blood clot formation, sparing small vessels (capillaries, venules) and not associated with ANCA, with an incidence of 2--3 per million, a slight male predominance, and affecting all ages. The condition involves immune complex deposition in the arterial wall (especially in HBV-associated PAN), complement activation, and neutrophil recruitment leading to vessel wall tissue death, with approximately 10--20\% of cases associated with chronic hepatitis B virus (HBV) infection.
\section{Causes}
\section{Risk Factors}

\begin{itemize}[noitemsep]
  \item Genetic predisposition and family history
  \item Advancing age
  \item Lifestyle factors (diet, exercise, smoking, alcohol)
  \item Environmental exposures
  \item Underlying medical conditions
  \item Medication use
\end{itemize}
\section{Signs and Symptoms}

\subsection{Common symptoms}
\begin{itemize}[noitemsep]
  \item You may experience constitutional symptoms (fever, weight loss, muscle pain, joint pain), mononeuritis multiplex (the most characteristic neurological manifestation), livedo reticularis, testicular pain, kidney artery involvement (kidney tissue death due to lack of blood supply, high blood pressure, but not glomerulonephritis---distinguishing PAN from microscopic polyangiitis), mesenteric reduced blood flow, and cutaneous manifestations (nodules, ulcers, gangrene). 
  \end{itemize}

\subsection{Emergency warning signs}
\begin{itemize}[noitemsep]
  \item Severe or worsening symptoms of polyarteritis nodosa require immediate medical evaluation
\end{itemize}
\section{Progression and Stages}
The condition may progress through different stages of severity over time. Early stages often involve mild or subtle symptoms that may be overlooked. Moderate stages involve more noticeable symptoms affecting daily function. Advanced stages may involve significant impairment and complications.
\section{Complications}
\begin{itemize}[noitemsep]
  \item Disease progression leading to worsening symptoms
  \item Secondary infections or injuries
  \item Side effects of treatment
  \item Reduced quality of life
  \item Emotional and psychological impact
\end{itemize}
\section{Diagnosis}
\begin{itemize}[noitemsep]
  \item To make the diagnosis, doctors look for biopsy of affected tissue showing necrotizing arteritis with fibrinoid tissue death, or angiography showing microaneurysms and irregular stenoses. Treatment for HBV-negative PAN is high-dose corticosteroids plus cyclophosphamide for severe disease
  \item HBV-associated PAN requires antiviral therapy plus a shorter course of corticosteroids and plasmapheresis.
  \item Comprehensive medical history and physical examination
  \item Laboratory tests to assess organ function
  \item Imaging studies to visualize affected structures
  \item Specialized testing depending on the affected system
  \item Consultation with relevant specialists
\end{itemize}
\section{Prevention}
\begin{itemize}[noitemsep]
  \item Maintain a healthy lifestyle with balanced diet and exercise
  \item Avoid known risk factors
  \item Attend regular medical check-ups for early detection
  \item Follow recommended screening guidelines
\end{itemize}
\section{Treatment Options}
Medications to manage symptoms and address underlying mechanisms Lifestyle modifications and self-care measures Physical therapy and rehabilitation Surgical intervention when indicated Regular monitoring and follow-up care Patient education and support services

\begin{itemize}[noitemsep]
  \item Medications to manage symptoms and address underlying mechanisms
  \item Lifestyle modifications and self-care measures
  \item Physical therapy and rehabilitation
  \item Surgical intervention when indicated
  \item Regular monitoring and follow-up care
\end{itemize}
\section{Living with It}
\begin{itemize}[noitemsep]
  \item Follow your treatment plan as prescribed
  \item Maintain regular follow-up with your healthcare provider
  \item Make necessary lifestyle adjustments
  \item Seek support from family, friends, or support groups
  \item Stay informed about your condition
\end{itemize}
\section{Outlook}
The outlook is good with treatment, but severe cases can be fatal. The outlook varies depending on the specific condition, severity at diagnosis, and response to treatment. Early intervention generally leads to better outcomes. Many conditions can be effectively managed with appropriate treatment. Regular follow-up care is important for monitoring and treatment adjustment.
